\documentclass[11pt]{article}
\usepackage[letterpaper,margin=1in]{geometry}
\usepackage{fontspec}
\usepackage{amsmath,amssymb,amsthm}
\usepackage{graphicx,booktabs,xcolor,microtype}
\usepackage[numbers,sort&compress]{natbib}
\usepackage[font=small,labelfont=bf]{caption}
\usepackage[hidelinks,pdfusetitle]{hyperref}
\usepackage{enumitem}
\setlist{itemsep=1pt,topsep=3pt}
\graphicspath{{figures/}}
\newtheorem{proposition}{Proposition}
\newcommand{\R}{\mathbb{R}}
\newcommand{\ShuffleSnap}{ShuffleSnap}
\newcommand{\panel}[2]{\par\begin{center}\begin{minipage}{\linewidth}\centering
  \includegraphics[width=\linewidth]{#1}\captionof{figure}{#2}\end{minipage}\end{center}}

\title{\ShuffleSnap: strided windowed exact reassignment\\for point-to-grid layout}
\author{Kyle McDonald}
\date{October 6, 2026}

\begin{document}
\maketitle

\begin{abstract}
Snapping each point of a 2D embedding, such as a UMAP layout, to its own grid cell while
minimizing the total squared displacement is a linear assignment problem too large for dense
solvers at $10^5$ points. \ShuffleSnap, a local search inspired by Shellsort, splits the grid
into interleaved subgrids at a decreasing sequence of strides and solves small exact
assignment problems in shifted windows of each, starting every solve from closed-form
optimal-transport duals. Against optima certified by a sparse exact solver that proves
optimality for the full, dense problem, the default preset is within 0.06\% of the optimum on
nine of ten distributions at $10^5$ points, in a median of 2.3\,s on one core or 0.19\,s on
24 threads. The fast and default presets take about as long as BSP-OT with 16 and 64 plans,
and at equal time they are closer to the optimum on eight and nine of the ten distributions;
RasterFairy is 17--425\% above the optimum on seven. The remaining error forms long loops of
one-cell shifts that no window can remove, which makes uniform points, concave masks and
surplus cells hard.
\end{abstract}

\section{Introduction}

A common way to show a large image or document collection is to embed the items in the
plane, for example with t-SNE~\citep{vandermaaten2008tsne} or UMAP~\citep{mcinnes2018umap},
and then snap every item to its own cell of a regular grid, so that thumbnails tile the
screen without overlap. Snapping is an assignment problem: each point must receive a
distinct cell, and points should move as little as possible. Measured as total squared
displacement, the best snapping is a linear assignment problem, and also a discrete
optimal-transport problem with quadratic cost~\citep{peyre2019computational}.

For a few thousand points, practitioners solve it exactly with a dense Hungarian or
Jonker--Volgenant solver~\citep{kuhn1955hungarian,jonker1987shortest}. The dense cost
matrix needs $8N^2$ bytes, 80\,GB at $N=100{,}000$, so these solvers stop being usable
between roughly 10{,}000 and 30{,}000 points. Beyond that, people use recursive splitting,
as RasterFairy does~\citep{klingemann2015rasterfairy}, or fast approximate transport such as
BSP-OT~\citep{genest2025bspot}.

\ShuffleSnap\ takes its name and its idea from Shellsort~\citep{shell1959highspeed}, which
sorts interleaved subsequences at a decreasing sequence of gaps. At stride $s$, \ShuffleSnap\
splits the grid into $s^2$ interleaved subgrids and solves small \emph{exact} assignment
problems inside $8\times8$ windows of each, dividing $s$ by $\sqrt2$ until it reaches 1
(Figure~\ref{fig:schematic}). All work is at full resolution; nothing is aggregated.

We measure every gap against a proven optimum and report where the method fails. We
contribute a careful parallel implementation with a 70-line serial reference that reproduces
it exactly; a window solver that starts from closed-form transport duals; a certified exact
solver that reaches $10^6$ points, so that every gap is measured against a proven optimum;
measured design choices; a diagnosis of the remaining error, with inputs on which the relative gap
is unbounded; and a benchmark against exact solvers, RasterFairy and BSP-OT that keeps
displacement and neighbourhood preservation apart.

\clearpage

\section{Problem}
\label{sec:problem}

A target is a boolean mask on a $W\times H$ lattice; its usable positions are the cells
$c_1,\dots,c_M$, centered on integer coordinates so that one unit is one cell. Input points
are mapped into this frame before anything else, each axis of their bounding box onto that
of the cells, and every method receives the same mapped points $p_1,\dots,p_N$ ($N\le M$).
We seek an injective assignment $\sigma$ minimizing
\begin{equation}
  C(\sigma) = \sum_{i=1}^N \lVert p_i - c_{\sigma(i)}\rVert^2 .
  \label{eq:cost}
\end{equation}
When $N<M$ the solver must also decide which cells stay empty, and the cell term of the
cost then depends on that choice. The linear-programming relaxation of
Equation~\eqref{eq:cost} has integral optima, because its constraint matrix is totally
unimodular, so linear-programming duality gives a certificate of optimality.

\begin{proposition}\label{prop:lb}
For any cell potentials $v\in\R^M$,
$\mathrm{LB}(v) = \sum_{i} \min_{j} (\lVert p_i-c_j\rVert^2-v_j) + \sum_{j} v_j - (M-N)\max_j v_j$
is a lower bound on the optimal cost, and it is tight for optimal potentials.
\end{proposition}
\begin{proof}[Sketch]
Shifting $v$ by its maximum gives nonpositive potentials $v'$; with $u_i=\min_j(c_{ij}-v'_j)$,
the pair $(u,v')$ is feasible for the dual of the assignment program, and its dual objective
equals $\mathrm{LB}(v)$. Weak duality gives the bound, strong duality its tightness.
\end{proof}

The bound needs a minimum over \emph{all} cells for every point, but it does not depend on
how $v$ was found, so $C(\sigma)-\mathrm{LB}(v)$ certifies the gap of any assignment. We
report the relative excess $C(\sigma)/C^\star-1$ over a certified optimum $C^\star$, and
where it misleads, the absolute excess per point: points that almost sit on the cells have a
tiny optimum, which inflates relative numbers, and heavy-tailed inputs have a huge one, which
hides them.

\panel{fig_schematic}{(a)~A window collects the points that occupy its cells and
reassigns them among the same cells by an exact linear assignment. (b)~At stride $s$ a
window takes every $s$-th row and column: it keeps 64 cells but spans $7s+1$ cells per
axis, so points can travel far. (c)~Each subgrid is tiled four times, shifted by half a
window, so that points can cross every window border.\label{fig:schematic}}

\clearpage

\section{\ShuffleSnap}
\label{sec:method}

\ShuffleSnap\ has two main parameters, the window side $w$ and the stride ratio $r$; the
default is $w=8$ and $r=\sqrt2$ (Section~\ref{sec:impl}).
\begin{enumerate}
\item \textbf{Start.} Place the points on a random subset of the cells, as a seeded
  random permutation.
\item \textbf{Strided windows.} At stride $(s_x,s_y)$, lattice positions with the same
  residues $(x \bmod s_x, y\bmod s_y)$ form one of $s_xs_y$ subgrids. Each subgrid is
  tiled with windows of $w\times w$ of its cells, spaced $s$ apart and spanning $(w-1)s+1$
  cells per axis.
\item \textbf{Exact reassignment.} In every window, the points that occupy its cells are
  reassigned among the same cells, empty ones included, by an exact linear assignment on
  squared distances to the original coordinates.
\item \textbf{Shifted tilings.} The tiling is repeated with shifts $(0,0)$, $(0,w/2)$,
  $(w/2,0)$ and $(w/2,w/2)$ in subgrid coordinates. These four phases form a round; the windows
  of a phase are disjoint and are solved in parallel.
\item \textbf{Schedule.} On an axis of length $L$ the first stride is the smallest power of
  two with $ws\ge L$. After one round both strides are divided by $r$ and rounded, never
  below 1. Stride-1 rounds repeat until a round moves no point, or until a time budget runs
  out.
\end{enumerate}

\paragraph{Acceptance.} Each window is a problem with $k$ occupants and $m\ge k$ cells,
which we solve exactly by shortest augmenting paths in $O(k^2m)$ time, from geometric starting
duals (Section~\ref{sec:impl}). The new assignment replaces the
old one only if it is cheaper by more than a relative $10^{-12}$. Since the old assignment
is feasible for the window, every accepted change lowers the global cost: the cost is
monotone, the search terminates, and ties, such as duplicated points, never cause cycling.
The final no-move round certifies a fixed point of the windows, not global optimality.

\paragraph{Boundaries, rectangles, masks.} Windows that cross a subgrid edge are clipped,
never dropped, on all four sides, so every lattice position belongs to exactly one window of
every phase and boundary cells trade with the interior; windows with fewer than two usable
cells are skipped. Each axis starts from its own stride, and the axis that reaches 1 first
stays there. Masked
positions never enter a window. Empty usable cells do, so with more cells than points the
method also decides which cells stay empty. A prime number of points can therefore fill an
exact-fit grid with a partial last row, or a full rectangle with a few surplus cells.

\panel{fig_progress}{Progress on 10{,}000 points of a UMAP embedding of MNIST
digits~\citep{lecun1998gradient}, default preset, every second stride shown. Every point
keeps the color of its input position, and the dark bands are graph-paper lines drawn on the input; both travel with the
points, so smooth colors and unbroken lines mean that the input neighbourhoods survived.
Numbers are the mean squared displacement in cells$^2$, next to the certified
optimum.\label{fig:progress}}

\clearpage

\subsection{Implementation and design changes}
\label{sec:impl}

The kernel is C++17 with nanobind bindings, parallelized with OpenMP over the windows of a
phase. Because windows in a phase are disjoint and each solve is deterministic, the result
is bit-identical for any thread count. A window whose cells have not changed occupant since
it was last solved in the same stage is already optimal and is skipped, which matters in the
stride-1 rounds. The exact cost is recorded after every phase, which gives an anytime trace.
A serial reference of about 70 lines of NumPy and SciPy reproduces the kernel's assignments
exactly and serves as an executable specification. A phase at any stride has about $M/w^2$
windows of at most $w^2$ cells, so each level costs about the same, and with $O(\log L)$
levels the schedule costs $O(M\log M)$ plus the convergence rounds; the working memory is
55--58 bytes per point.

We tuned only on development instances and report every claim on held-out ones.
\emph{Window size} is the main quality knob: each step from $8\times8$ to $12\times12$ to
$16\times16$ cost 2--2.6 times the time and removed one half to two thirds of the remaining
excess. The \emph{fast} ($5\times5$) and \emph{balanced} ($8\times8$, default) presets take
about as long as BSP-OT with 16 and 64 plans; \emph{quality} is $12\times12$. \emph{Stride ratio:}
dividing by $\sqrt2$ instead of 2 was the best or close to the best ratio at every window
size. \emph{Start:} structured starts, such as
recursive bisection or sorting by rows, are worse than a random one: they create straight
seams that windows cannot repair. \emph{Solver:} every point of a distant cluster prefers
the same corner cell, which defeats the usual greedy start. Starting from the Brenier
potential~\citep{brenier1991polar} of the per-axis sorted map instead stays exact and made
window solves 2--5 times faster, which buys larger windows at equal time;
Jonker--Volgenant and an auction were slower.

\section{A certified exact solver}
\label{sec:exact}

At $10^5$ points no dense matrix fits in memory, so we built a solver that never forms the
dense problem, yet certifies optimality against it. It gives each point candidate edges:
its 12 nearest cells, plus the $5\times5$ lattice neighbourhood of its cell in a starting
assignment, which guarantees that a perfect matching exists. It solves the min-cost flow on
these edges with OR-Tools' cost-scaling solver~\citep{ortools}, on costs scaled to integers.
From that solution it computes cell potentials by shortest paths, in exact integer
arithmetic, and evaluates $\mathrm{LB}$ (Proposition~\ref{prop:lb}) over \emph{all} cells,
scanning blocks of cells only where they can beat the current minimum. The pass also lists
every cell whose reduced cost undercuts a point's current cell; if any exists, those edges
are added and the flow is solved again. This is the core-and-price idea of sparse
assignment and transport solvers~\citep{volgenant1996linear,schmitzer2016sparse}. The
certificate $C(\sigma)-\mathrm{LB}(v)$ is computed from real-valued costs and holds whatever
errors occur along the way; integer rounding can only loosen it, by at most $4N/S$ for scale
$S$. The result is therefore the exact optimum of the dense problem. Wherever SciPy's dense
solver fits, both found the same optimum: on 600 small instances and on all ten
distributions at $10^4$ points. Started from
\ShuffleSnap's layout, it certified all 202 benchmark instances up to $2\times10^5$ points,
and uniform instances up to $10^6$, in 1 to 13 pricing rounds. Started from bisection
instead, it failed on four of ten instances at $10^5$ points.

\paragraph{Correctness checks.} We checked three properties independently of the kernel.
Every output of every method gave each point its own usable cell (all 1{,}217 completed runs of
the main benchmark); the exact cost recorded after every phase was monotone in all 540
\ShuffleSnap\ runs and ended at the independently recomputed cost; and the window solver
matched brute-force enumeration and SciPy on random, tie-heavy and badly scaled matrices.

\clearpage

\section{Experimental setup}
\label{sec:setup}

All runs used one workstation with an Intel i9-12900KS (8 performance and 8 efficiency
cores, 24 hardware threads) and 62\,GB of memory. Single-thread runs were pinned to one
performance core; multi-thread runs used all 24 threads. Every run was a fresh process after
a warm-up call on a small instance. Its time includes all of the method's own work, such as
cost matrices, trees, starts and internal normalization, but not generating the instance or
checking the result. Memory is the rise of the peak resident set during the call. Design
decisions used instance seed 0, and all results use held-out seeds 1--3.

The ten synthetic distributions are uniform, Gaussian, a 10-component anisotropic
mixture like typical embeddings, two clusters with 90/10 mass, a thin ring, a noisy spiral,
a checkerboard density, heavy tails, 16 heavily duplicated sites and a nearly degenerate line.
The real data are UMAP embeddings~\citep{mcinnes2018umap} of all 70{,}000 MNIST and
Fashion-MNIST~\citep{xiao2017fashion} images. The targets are near-square exact-fit grids,
rectangles with surplus cells, a 4:1 rectangle, a disk, a ring and a blocky, rectilinear
S-shaped mask.

Each comparator received the same coordinates and cells, and every output was checked for
validity independently. \emph{RasterFairy}'s 2019 PyPI release orders points by
integer-truncated coordinates, which discards sub-cell order unless inputs are scaled up; its
current upstream version normalizes internally and solves leaf blocks of up to 200 points
exactly, and is our representative setting. \emph{BSP-OT} ran through its official Python
binding with 16 and 64 merged plans. The binding exposes no seed, so repeated calls differ:
the costs of the same instance with 1 and 24 threads differed by $-8\%$ to $+13\%$, with no
trend. We measure neighbourhood preservation separately, as the fraction of each point's 10
nearest input neighbours among its 10 nearest grid neighbours (kNN recall), and the fraction
of a point's occupied neighbour cells that hold one of its 8 nearest input neighbours
(lattice precision).

\begin{center}
\begin{minipage}{\linewidth}\centering
\captionof{table}{Excess over the certified optimum (\%) at $N=100{,}000$, mean of three
held-out instances, one thread. The last column is the certified solver's time, started
from \ShuffleSnap\ (instance seed 1); the bottom row is the median time of each
method.\label{tab:main}}
\small
\input{tables/tab_main_compact_100k_t1}
\end{minipage}
\end{center}

\clearpage

\section{Results}
\label{sec:results}

\subsection{Quality and time at $10^5$ points}

\panel{fig_tradeoff_100k}{Excess over the certified optimum against single-thread time at
$N=100{,}000$ (mean over three held-out instances; bars span their range). The three
\ShuffleSnap\ presets form the blue curve, from fast ($5\times5$) to quality ($12\times12$). The black line is the
time of the certified exact solver. Heavy tails are omitted: every method is below 0.001\%
there.\label{fig:tradeoff}}

On nine of the ten distributions the default preset is within 0.06\% of the optimum, and
on six within 0.005\% (Table~\ref{tab:main}, Figure~\ref{fig:tradeoff}). It needs a median of
2.3\,s on one core. The exception is uniform points (0.84\%). The presets trade time for
quality on every distribution: fast (median 0.77\,s) is within 0.16\% on the same nine, and
quality (6.7\,s) lowers the uniform excess to 0.44\%. At equal time \ShuffleSnap\ beats BSP-OT:
fast is closer to the optimum than BSP-OT with 16 plans (0.71\,s) on eight of ten
distributions, and the default than 64 plans (2.9\,s) on nine, e.g.\ 2.6\% and 0.84\% against
45\% and 32\% on uniform points. BSP-OT is ahead only on the near-line and the 90/10 clusters.
RasterFairy takes about 3\,s and is 17--425\% above the optimum on seven of the ten
distributions; it does well on the Gaussian and the 90/10 clusters. With 24 threads, \ShuffleSnap\
ran 7--14 times faster with bit-identical layouts (0.19\,s for the default, 0.09\,s for
fast), while BSP-OT gained 4--7 times, because the cycle merge that combines its plans is largely sequential.

\clearpage

\subsection{Matched budgets and scaling}

\panel{fig_anytime_t1}{Excess over the certified optimum against elapsed time, one thread.
\ShuffleSnap's step curves record the exact cost after every phase, so each shows what a
time budget ending at that moment would return. BSP-OT runs with 1 to 1024 plans;
RasterFairy with exact leaf blocks of 1 to 1000 points.\label{fig:anytime}}

BSP-OT improves slowly with more plans and levels off far from the optimum: from 1 to 1024
plans, its excess drops from 249\% to 28\% on uniform points and from 16\% to 1.6\% on MNIST
(Figure~\ref{fig:anytime}). The fast \ShuffleSnap\ preset passes BSP-OT's best result after
0.2--0.5\,s and the default after 0.5--1.3\,s.
With less than 0.12--0.36\,s on one thread, BSP-OT is the better choice, because \ShuffleSnap\
has only finished its coarse levels; with 24 threads, the fast preset is ahead at every budget.

\panel{fig_scaling}{Time against $N$ on one thread. Dense SciPy ran up to 50k uniform and
20k mixture points; the certified solver up to $10^6$ uniform and $2\times10^5$ mixture
points.\label{fig:scaling}}

\ShuffleSnap, BSP-OT and RasterFairy all scale nearly linearly (Figure~\ref{fig:scaling}):
$10^6$ uniform points take the fast preset 9.8\,s on one core, as long as BSP-OT with 16
plans, and the default 21\,s (1.5\,s on 24 threads) with 55 bytes per point, a quarter to a twelfth of
the comparators' memory. The certified solver took 11\,min and 4.2\,GiB there; dense SciPy
needed 6\,min and 23\,GiB at 50{,}000. The default's uniform excess grows with $N$: 0.1\% at $10^3$, 0.84\% at $10^5$, 1.3\% at $10^6$.

\clearpage

\subsection{Real embeddings: displacement versus neighbourhoods}

\panel{fig_gallery_mnist}{All 70{,}000 MNIST digits, UMAP layout snapped to a
$265\times265$ grid. Colors and graph-paper lines follow each point's input position, so
bent lines show how a method deforms the layout, and broken ones where it tears it; numbers
are the mean squared displacement. Row~2 enlarges the white box, row~3 shows each point's
distance to its cell in the optimum (log scale; white is the optimal cell), row~4 its lattice
precision. \ShuffleSnap\ leaves 74\% of the points on their optimal
cell; BSP-OT's median point is 8 cells away.\label{fig:visual}}

On both embeddings the default preset is within 0.002\% of the certified optimum, the fast
preset within 0.031\% (BSP-OT with 16 plans: 3.7--9.2\%), and their layouts match the optimal one
(Figure~\ref{fig:visual}). The two measures disagree on BSP-OT: with 64 plans its
displacement is only 2.4\% (MNIST) and 6.7\% (Fashion-MNIST) above the optimum, yet kNN
recall falls to 0.31--0.33 against 0.61--0.65, and lattice precision to about half. Its
merged partitions are close to the transport map on average but shuffle points locally.
RasterFairy shows the opposite: its displacement is 20--25\% above the optimum, yet its
neighbourhood measures are close to the optimum's, because local order survives inside each
slice. Even the optimum tears the layout along cluster boundaries (pale lines in row~4).
Minimizing displacement and preserving neighbourhoods are different goals.

\clearpage

\subsection{Masks, surplus cells and design choices}

\panel{fig_masks_row}{20{,}000 points of the MNIST embedding on a rectilinear S-shaped mask and on a
rectangle with 10\% surplus cells; numbers are the excess over the certified optimum. With
surplus cells, the optimum leaves one region empty, far from the data. BSP-OT's partial
matching leaves ragged holes; RasterFairy removes surplus cells at random before
slicing.\label{fig:masks}}

On disks and rings, the default stays within 0.07\% of the optimum (fast 0.14\%), while
BSP-OT is 2.5--9.3\% and RasterFairy 11--54\% above it. On the S-shaped mask, coarse
schedules can route mass across the bays. With uniform points the default reaches 1.4\%,
against 9.3\% for BSP-OT with 64 plans, and quality 0.88\%; this is the one mask input where
the fast preset loses to BSP-OT with 16 plans (17\% against 12\%). Surplus cells, where the solver must choose
the empty cells, are harder: the default stays within 0.91\% with 10\% surplus but reaches 3.1\%
on uniform points with less than a row of surplus. The comparators suffer most here: with
10\% surplus, BSP-OT's partial matching is 18--109\% above the optimum and RasterFairy
72--350\%. For a prime
number of points, an exact-fit grid with a partial last row is therefore simpler and better
than a full rectangle for every method.

Varying one design choice at a time on held-out instances confirmed the development
findings. The window side is the main quality knob: on uniform points the excess falls from
3.8\% for $4\times4$ windows to 0.09\% for $24\times24$, while the time grows 28-fold.
Coarse levels are essential: capping the first stride made results up to 90 times worse
and up to 45 times slower, because stride-1 windows alone stall far from the optimum. So are
shifted tilings: with a single tiling no point can cross a window border (104\% on uniform
points), while two tilings give 2.0\%, four 0.84\% and nine 0.70\% at twice the time.
Skipping unchanged windows saves 22--24\% of the time without changing any result; polishing
at stride 2 after convergence changed almost nothing, and merging independent runs helped less
per second than a larger window.

\clearpage

\subsection{Anatomy of the remaining error}

\panel{fig_failures_row}{(a,\,b)~Where \ShuffleSnap\ and the optimum disagree on uniform
points, the difference is a few long closed loops of one-cell moves (loops longer than 30
points in blue); at $10^5$ points loops of more than 300 points carry most of the excess.
(c,\,d)~Points drawn inside the S mask (optimum $\approx$ identity): $6\times6$ windows with
halving strides route points across the bays, where they line up; the default leaves fewer.
Numbers are mean squared displacements.\label{fig:failures}}

Why does a method that solves every window exactly miss the optimum? The permutation that
turns a \ShuffleSnap\ layout into the optimal one splits into cycles, and almost all of the
excess sits in a few very long ones, in which every point is shifted by about one cell
(Figure~\ref{fig:failures}a,b). Inside a window, part of such a loop is a run of consistently
shifted points, which is already optimal for that occupant set, and a strided window cannot
move a point by one cell. On points placed exactly on the cells, with an axis-aligned ring
of cells rotated by one step, the full schedule never moves a ring larger than the window,
so the relative excess over the optimum of zero is unbounded. The same mechanism appears at
a larger scale on masks: coarse windows route mass across concave bays, which no stride-1
window can undo (Figure~\ref{fig:failures}c,d), and a mask can even cut two usable cells off
from each other when no scheduled window contains both.

\section{Discussion}

Coarse-to-fine strided windows get very close to the optimum on the smooth, clustered
layouts that motivate the problem, in a fraction of a second with all cores at $10^5$ points,
deterministically and with little memory. They fail where a local search must: on uniform
data, lattices and surplus cells, whose errors need moves no window can see.
A billion points converged in 1\,h\,43\,min using 53\,GiB, passing a six-hour run of an
earlier implementation after 44 minutes; a single-plan BSP-OT ended 8.2\% above.
Built as an evaluation tool, the exact solver proved optimality at
$10^5$ points in a median of about two minutes when started from \ShuffleSnap. When a proven
optimum is worth that wait, it is the better choice; otherwise the default preset is the
fast, near-optimal answer, and BSP-OT is best only under budgets below 0.12--0.36\,s on one
core.
If neighbourhoods matter, they should be optimized or measured directly, for example with
LAS/FLAS~\citep{barthel2023improved}.

\paragraph{Limitations.} Everything ran on one workstation, CPU only (BSP-OT's GPU version
was not tested), for squared Euclidean displacement in 2D, on two real UMAP embeddings, with
representative settings of each comparator. The package is on PyPI as
\texttt{shufflesnap}; raw results and figure scripts are in a separate repository.
\label{page:body-end}

\clearpage
\bibliographystyle{plainnat}
\bibliography{refs}
\end{document}
